% \iffalse meta-comment
%
%% regulatory-html.dtx
%% Copyright 2024-2026 E. Nijenhuis
%
% This work may be distributed and/or modified under the
% conditions of the LaTeX Project Public License, either version 1.3c
% of this license or (at your option) any later version.
% The latest version of this license is in
% http://www.latex-project.org/lppl.txt
% and version 1.3c or later is part of all distributions of LaTeX
% version 2005/12/01 or later.
%
% This work has the LPPL maintenance status ‘maintained’.
%
% The Current Maintainer of this work is E. Nijenhuis.
%
% This work consists of the files listed in the meta-comment of
% regulatory-struct.dtx.
%
% \fi
%
% \iffalse
%<*driver>
\ProvidesFile{regulatory-html.dtx}
%</driver>
%<html>\ProvidesFile{regulatory.4ht}
%<html>    [2026/09/05 0.0.5 Xerdi's Regulatory Package (tex4ht configuration)]
%
%<*driver>
\documentclass[10pt,english]{ltxdoc}
%! suppress = InclusionLoop
\usepackage{regulatory}
\usepackage{tabularx}
\usepackage[english,dutch]{babel}
%% regulatory-preamble.tex
%% Copyright 2024 E. Nijenhuis
%
% This work may be distributed and/or modified under the
% conditions of the LaTeX Project Public License, either version 1.3c
% of this license or (at your option) any later version.
% The latest version of this license is in
% http://www.latex-project.org/lppl.txt
% and version 1.3c or later is part of all distributions of LaTeX
% version 2005/12/01 or later.
%
% This work has the LPPL maintenance status ‘maintained’.
%
% The Current Maintainer of this work is E. Nijenhuis.
%
% This work consists of the files regulatory.ins,
% regulatory-struct.dtx, regulatory-ref.dtx,
% regulatory-defs.dtx, regulatory-attachments.dtx,
% regulatory-md.dtx,
% regulatory-html.dtx,
% regulatory-english.def, regulatory-dutch.def,
% regulatory.tex, regulatory-preamble.tex,
% regulatory-nl.tex, regulatory-en.tex,
% example1.bib, example2.bib,
% example1-nl.tex, example2-nl.tex,
% example1-en.tex, example2-en.tex,
% md-example.tex, example.md,
% and the derived files regulatory.sty, regulatory-struct.sty,
% regulatory-defs.sty, regulatory-ref.sty,
% regulatory-attachments.sty, regulatory-md.sty
% and regulatory.4ht.
\usepackage{gitinfo-lua}
\usepackage{listings}
\usepackage{adjustbox}
\usepackage{tikz}
\usetikzlibrary{arrows.meta}

\usepackage{fontspec}
\usepackage{textcomp}

\def\projecturl{https://github.com/Xerdi/regulatory}

\definecolor{greycolor}{HTML}{BFBFBF}
\definecolor{primary}{HTML}{174A67}
\definecolor{darkcolor}{HTML}{091D28}
\definecolor{greencolor}{HTML}{176734}
\colorlet{artlabel}{primary}
\colorlet{green}{greencolor}

\lstdefinestyle{tex}{%
    language=[LaTeX]TeX,
    keywordsprefix={\\},
    alsoletter={\\},
    morekeywords={article,para,textfill,masterdocument,refdocument,loadglsdefs,printdefs,setmiddleconjunction,setlastconjunction,setrangeconjunction,setconjunction,markdownInput,glssetwidest,describe}
}

\lstdefinestyle{bash}{%
    language=bash,
    morekeywords={pdflatex,lualatex,latexmk,makeglossaries,bib2gls}
}

\lstdefinelanguage{Markdown}[LaTeX]{TeX}{%
    sensitive=false,
    morecomment=[l][\bfseries\ttfamily\color{green}]{\#},
    morecomment=[s][\bfseries\ttfamily\color{purple}]{:\ \ }{\ },
    morestring=[s][\color{black}]{\{}{\}},
    morekeywords={article,para,textfill,masterdocument,refdocument,loadglsdefs,printdefs,setmiddleconjunction,setlastconjunction,setrangeconjunction,setconjunction,markdownInput,glssetwidest,describe},
    keywordsprefix={\\},
    alsoletter={\\},
}

\lstdefinestyle{md}{%
    language=Markdown
}

\lstdefinelanguage{BibTeX}{%
    keywords={%
        article,book,collectedbook,conference,electronic,entry,ieeetranbstctl,%
        inbook,incollectedbook,incollection,injournal,inproceedings,%
        manual,mastersthesis,misc,patent,periodical,phdthesis,preamble,%
        proceedings,standard,string,techreport,unpublished%
    },
    keywordsprefix={@},
    comment=[l][\itshape]{@comment},
    sensitive=false,
}

\lstdefinestyle{bib}{%
    language=BibTeX,
    morestring=[s][\bfseries\ttfamily\color{purple}]{\ \ \ \ }{\ },
    morecomment=[n][\bfseries\ttfamily\color{green}]{\ \{}{\}},
}

\lstset{
    title=\lstname,
    basicstyle=\ttfamily,
    numberstyle=\ttfamily,
    numbersep=5pt,
    rulecolor=\color{greycolor},
    stringstyle=\color{pink},
    keywordstyle=\color{primary},
    commentstyle=\itshape\color{darkcolor},
    breaklines=true,
    captionpos=b,
    tabsize=4,
    showtabs=true
}

\newcommand\package{\texttt}
\newcommand\file{\texttt}
\newcommand\option{\texttt}

% Reads test/conformance.tex, which test/conformance.sh generates and which is
% kept in the repository: veraPDF is not available everywhere this manual is
% built, so the verdicts travel with the sources.
\newcommand\conformanceversion{?}
\newcommand\conformanceimage{?}
% A digest carries no break points of its own and does not fit a line, so this
% offers one after every character. The argument is expanded first, otherwise the
% loop sees one macro instead of the characters it stands for.
\newcommand\anywherebreakend{}
\def\anywherebreakloop#1{%
    \ifx#1\anywherebreakend\else#1\allowbreak\expandafter\anywherebreakloop\fi}
\newcommand\anywherebreak[1]{\anywherebreakloop#1\anywherebreakend}
\newcommand\conformanceimagetext{%
    \expandafter\anywherebreak\expandafter{\conformanceimage}}

% \ignorespaces, because this macro typesets nothing while its line in the
% generated file still ends in a newline, and that space would otherwise open
% the first cell of the table.
\newcommand\conformancevalidator[2]{%
    \gdef\conformanceversion{#1}\gdef\conformanceimage{#2}\ignorespaces}
\newcommand\conformancerow[6]{%
    \texttt{#1} & \texttt{#2} & \texttt{#3} & \texttt{#4} & \texttt{#5}%
    \ifthenelse{\equal{#6}{}}{}{ \footnotesize(#6)} \\}

% Points at the resulting PDF of an example, which is attached to this manual.
% Its argument is the label of the artifact the example was declared with, and it
% falls back to plain text whenever the PDF isn't there.
\newcommand\exampleresult[1]{%
    \par\noindent
    \translation{The result of this example is}{Het resultaat van dit voorbeeld is}
    \documentattachment{#1}{\file{\artifactfilename{#1}}}.\par}

\def\packagedate{\gitdate}
\def\packageversion{\gitversion}

%\def\cmd{\lstinline[style=tex]}
\def\cmditem#1{\item[\ttfamily\textbackslash{}#1]}

\newcommand\Vtextvisiblespace[1][.3em]{%
  \mbox{\kern.06em\vrule height.3ex}%
  \vbox{\hrule width#1}%
  \hbox{\vrule height.3ex}}

% Package zref-clever loads its language files lazily, which would otherwise
% happen inside a \DocInput region, where the percent sign is ignored and the
% comments of those files end up being parsed as keys. Touching every language
% of this manual schedules them at the start of the document instead, while the
% percent sign still is a comment character.
\zcsetup{lang=english}
\zcsetup{lang=dutch}
\zcsetup{lang=current}

\newcommand\translation[2]{#1}
\begin{document}
    \selectlanguage{english}
    \DocInput{regulatory-html.dtx}
\end{document}
%</driver>
% \fi
%
% \subsection{\texorpdfstring{\file{regulatory.4ht}}{regulatory.4ht}}
% \setcounter{CodelineNo}{0}
%
% This is the \package{tex4ht} configuration of \zcref{sec:html}, which is a source of its own
% rather than a module of the bundle: it is not a package, it is read by \package{tex4ht} and by
% nothing else, and it only exists at all because an HTML run replaces the machinery the printed
% headings are built on. It identifies itself with \cmd{\ProvidesFile} all the same, so that a
% document can tell whether it was found: \package{tex4ht} ships no configuration for this package,
% and a conversion without one fails silently rather than loudly (\zcref{sec:html}).
%
%
% \subsubsection{What this fixes}
%
%
% This file is what emits the structure of a regulatory document in HTML. It
% replaces \cmd{\article} and \cmd{\para} outright, so it does that in either of the two
% heading branches the package can take (see below), and without it there is no
% structure at all: an article becomes a paragraph with a bold run in it, and
% the output carries no \texttt{<h1>}..\texttt{<h6>} and no \texttt{<section>}. Measured on example1-nl,
% four articles:
%
% \begin{center}
%     \begin{tabular}{@{}lrrl@{}}
%         & \texttt{<section>} & \texttt{<h\meta{n}>} & \texttt{make4ht} \\
%         \hline
%         without this file & 0 & 0 & exit 0 \\
%         with it           & 5 & 5 & exit 0 \\
%     \end{tabular}
% \end{center}
%
% Note both exit codes. The conversion \emph{succeeds} without this file and produces a
% document with no headings; what falls away does not announce itself.
%
% A document that does not declare \cmd{\DocumentMetadata} hits a hard failure on top
% of that, because regulatory then builds \cmd{\article} with titlesec's \cmd{\titleclass}
% and tex4ht replaces the sectioning machinery titlesec hooks into, leaving
% \cmd{\titleformat} nothing to attach to. Measured on the same document with neither
% this file nor \cmd{\DocumentMetadata}: make4ht exits 1 on `No format for this
% command', `Missing number' and `Illegal unit of measure', and four faults arrive at
% once, none of which announces itself:
%
% \begin{enumerate}
%     \item \package{titlesec} errors per heading and leaks its spacing argument into the
%         text: `0ptEerste artikel'. Measured: five occurrences of `0pt' in the output.
%     \item It also consumes the first letter of the following text as an argument, so
%         `Een genummerde paragraaf' comes out as `0ptE' plus `en genummerde paragraaf'.
%         The text is not misformatted, it is broken apart.
%     \item Article numbers vanish. In a legal instrument the number is not decoration; it
%         is the address by which a clause is cited.
%     \item \cmd{\aref} still emits \verb|<a href="#article.1">|, but nothing emits an
%         element with that id. Measured: three such links, zero matching ids.
% \end{enumerate}
%
% That branch is the narrower case, not the reason this file exists. A document
% that attaches anything declares \cmd{\DocumentMetadata} already and never reaches
% titlesec -- and still needs this file for everything above.
%
% The \texttt{paras} environment needs nothing here: it is an ordinary enumitem list
% and tex4ht already renders it with correct numbering.
%
%
% \subsubsection{Two traps}
%
%
% This file must not use \cmd{\makeatletter} or \cmd{\makeatother}. The at sign is already a
% letter when the file runs (measured: \verb|\catcode`\@| is 11), so \cmd{\makeatletter} is
% redundant and the matching
% \cmd{\makeatother} is destructive: it turns @ back into an ordinary character for
% the rest of the package, so every \verb|\rref@|\dots{} internal silently becomes garbage.
% The symptom is a build that hangs and then dies inside color.4ht with
% \verb|Missing \begin{document}| -- an error naming a file that has nothing to do
% with the cause.
%
% (An earlier version of this comment explained the catcode by saying tex4ht
% inputs a .4ht `the moment it sees \cmd{\ProvidesPackage}'. That reason is wrong;
% only the conclusion held. See the next paragraph for what was measured.)
%
% The redefinitions are made directly, not deferred. Measured 2026-09-05
% against the split package (regulatory + -struct + -ref + -attachments), with
% a stub .4ht per package reporting its own state: tex4ht defers each file past
% the whole \texttt{usepackage} chain, not merely past the package that names it --
% \file{regulatory-struct.4ht} already sees \cmd{\attachdocument}, which regulatory-attachments
% defines afterwards. So \cmd{\article} and \cmd{\para} exist by the time we get here and we
% are overwriting them. \cmd{\AtEndOfPackage} looks like the careful choice and is not:
% it never fires, because by then no package is being loaded, and the
% redefinitions then silently never happen at all.
%
% Ordering caveat if this file is ever split along the package lines: tex4ht's
% own \file{titlesec.4ht} runs between the clusters: \file{regulatory.4ht} and
% regulatory-struct.4ht before it, regulatory-ref.4ht and
% regulatory-attachments.4ht after. An \cmd{\article}/\cmd{\para} fix must land in the first
% two. Keeping everything in one file keeps it on the safe side of that line.
%
%
% \subsubsection{Which heading branch this patches}
%
%
% regulatory-struct builds \cmd{\article} through titlesec's \cmd{\titleclass} only when
% \cmd{\ParseLaTeXeHeading} is absent. A bare \cmd{\DocumentMetadata}{} -- no testphase
% needed -- loads latex-lab-testphase-sec-template.sty and defines it, so such a
% document takes the \cmd{\DeclareInstance}{heading} branch and never reaches titlesec.
% Measured under both pdflatex and lualatex, 2026-09-05.
%
% That does not make this file optional for those documents. It replaces \cmd{\article}
% outright, so it is what emits the markup in either branch. Measured on one
% document with two articles, \cmd{\DocumentMetadata}{} on line 1, differing only in
% whether this file was present:
%
% \begin{center}
%     \begin{tabular}{@{}lrrl@{}}
%         & \texttt{<h\meta{n}>} & \texttt{regulatory-*} & \texttt{make4ht} \\
%         \hline
%         without this file & 0 & 0 & exit 0 \\
%         with it           & 2 & 6 & exit 0 \\
%     \end{tabular}
% \end{center}
%
% The id='article.N' anchors are hyperref's and appear either way. So without
% this file the conversion \emph{succeeds} and produces a document with no headings at
% all -- the same shape as the titlesec finding on the PDF side, where a clean
% build was no evidence either. What falls away does not announce itself.
% \iffalse
%<*html>
% \fi
%    \begin{macrocode}

\immediate\write-1{regulatory.4ht 2026-09-05}

\NewConfigure{regarticle}{4}
\NewConfigure{regpara}{4}
%    \end{macrocode}
% \iffalse
%</html>
% \fi
%
% The anchor id must be the one \cmd{\aref} links to. hyperref's \cmd{\refstepcounter} sets
% \verb|\@currentHref| to `article.N' -- exactly the `\#article.1' the broken output was
% already pointing at -- so emitting that as the element id is what turns those
% dead links into live ones.
%
% \iffalse
%<*html>
% \fi
%    \begin{macrocode}
\newcommand\reg@anchorid[1]{\HCode{ id="#1"}}
%    \end{macrocode}
% \iffalse
%</html>
% \fi
%
% Close any open paragraph before block-level markup, or the \texttt{<section>} is
% emitted inside a \texttt{<p>}. \cmd{\IgnorePar} stops tex4ht opening a fresh one immediately.
%
% \iffalse
%<*html>
% \fi
%    \begin{macrocode}
\newcommand\reg@blockmode{\ifvmode\IgnorePar\fi\EndP}
%    \end{macrocode}
% \iffalse
%</html>
% \fi
%
% A \texttt{<section>} opened by \cmd{\article} has to be closed by the next \cmd{\article} or at the
% end of the document. These are sectioning commands, not environments -- there
% is no \cmd{\endarticle} to hook -- so an explicit flag tracks it.
%
% \iffalse
%<*html>
% \fi
%    \begin{macrocode}
\newif\ifreg@inarticle \reg@inarticlefalse
\newif\ifreg@inpara    \reg@inparafalse

\newcommand\reg@closepara{%
  \ifreg@inpara \reg@blockmode\HCode{</section>}\reg@inparafalse \fi
}
\newcommand\reg@closearticle{%
  \reg@closepara
  \ifreg@inarticle \reg@blockmode\HCode{</section>}\reg@inarticlefalse \fi
}

\def\article{\@ifstar\reg@article@star\reg@article@plain}%
  %
  % titlesec's sectioning signature is \article*[toc-entry]{title}; both forms
  % are accepted so no existing document needs changing.
\newcommand\reg@article@plain[2][]{%
    \reg@closearticle
    \reg@blockmode
    \refstepcounter{article}%
    % \@currentlabel is what \label captures. Keep it the bare number so \aref
    % and \ref print the same thing the heading does.
    \def\@currentlabel{\thearticle}%
    \a:regarticle\reg@anchorid{\@currentHref}\HCode{>}%
    \b:regarticle\GetTranslation{Article}\ \thearticle.\c:regarticle
    #2%
    \d:regarticle
    \reg@inarticletrue
    \par\ShowPar
}
%    \end{macrocode}
% \iffalse
%</html>
% \fi
%
% Starred: an unnumbered heading, so the counter is deliberately not stepped
%
% \iffalse
%<*html>
% \fi
%    \begin{macrocode}
  % and no anchor is emitted -- there is no number for anything to cite.
\newcommand\reg@article@star[2][]{%
    \reg@closearticle
    \reg@blockmode
    \a:regarticle\HCode{>}%
    \b:regarticle\c:regarticle
    #2%
    \d:regarticle
    \reg@inarticletrue
    \par\ShowPar
}
%    \end{macrocode}
% \iffalse
%</html>
% \fi
%
% \subsubsection{The \cmd{\para} heading}
%
%
% \cmd{\para} is used, despite appearances. Grepping a legal source for \verb|\para| finds
% mostly \verb|\param|, and grepping for \verb|^\para| finds nothing because the calls are
% indented. It is what a document uses for sub-headings within an article, and
% leaving it out leaves those broken in a document whose articles all look
% correct.
%
% Nested inside the article's \texttt{<section>} rather than beside it: the sub-heading
% belongs to the article above it, and a flat structure would tell anything
% reading the outline otherwise.
% \iffalse
%<*html>
% \fi
%    \begin{macrocode}

\def\para{\@ifstar\reg@para@star\reg@para@plain}

\newcommand\reg@para@plain[2][]{%
  \reg@closepara
  \reg@blockmode
  \refstepcounter{para}%
  \def\@currentlabel{\thearticle.\thepara}%
  \a:regpara\reg@anchorid{\@currentHref}\HCode{>}%
  \b:regpara\thearticle.\thepara.\c:regpara
  #2%
  \d:regpara
  \reg@inparatrue
  \par\ShowPar
}

\newcommand\reg@para@star[2][]{%
  \reg@closepara
  \reg@blockmode
  \a:regpara\HCode{>}%
  \b:regpara\c:regpara
  #2%
  \d:regpara
  \reg@inparatrue
  \par\ShowPar
}
%    \end{macrocode}
% \iffalse
%</html>
% \fi
%
% \subsubsection{Cross-document references}
%
%
% A reference to another document must not become a hyperlink in HTML.
%
% In the PDF a link to another document is right whenever the reader holds both:
% it opens a document they were sent. On the web that no longer follows. A
% document cited with \cmd{\refdocument} or \cmd{\masterdocument} is commonly one issued to a
% counterparty rather than published, and a link to it would point at a file that
% is not there. A dead link in a legal text is worse than no link, since it
% suggests the reader is being denied something.
%
% This is a choice, not a limitation. A publisher who does put every cited
% document on the web wants the opposite, and gets it by redefining
% \verb|\rref@linkpr| after this file is read.
%
% The reference itself stays: the sentence still says which article of which
% agreement it means, which is what a reader needs in order to look it up in the
% copy they were sent. Only the anchor goes.
%
% regulatory already draws the distinction we need. \verb|\rref@current@url| is set by
% \verb|\rref@setlink| and is empty for a reference inside this document, non-empty
% for one resolved through zref-xr from another document's .aux. One test, at the
% single point where the package turns a reference into a link.
%
% Note this is a decision about the \emph{output format}, which is why it belongs here
% and not in the source. Writing \cmd{\aref}* in the .tex would drop the link from the
% PDF as well, where it is wanted.
% A reference whose anchor this format cannot produce must not look like a link
% either. In HTML a \texttt{paras} item carries no anchor under the name \cmd{\aref} points
% at (see the limitation above), so linking it would promise a jump that does
% not happen. The wording is what carries the meaning -- "het in eerste lid
% genoemde bedrag" tells the reader exactly where to look -- and a link that
% silently does nothing subtracts from that rather than adding to it.
%
% The counter is the discriminator, and it separates the two cases cleanly:
% \texttt{para} is the sectioning command, which this file \emph{does} anchor, while \texttt{parasi}
% and \texttt{parasii} are the enumerate levels of the \texttt{paras} list, which it cannot.
% Testing the counter rather than the type keeps a linked sub-heading linked.
% \iffalse
%<*html>
% \fi
%    \begin{macrocode}
\def\reg@parasictr{parasi}
\def\reg@parasiictr{parasii}
\newif\ifreg@anchorless
\newcommand\reg@checkanchorless{%
  \reg@anchorlessfalse
  \edef\reg@thisctr{\rref@current@counter}%
  \ifx\reg@thisctr\reg@parasictr \reg@anchorlesstrue \fi
  \ifx\reg@thisctr\reg@parasiictr \reg@anchorlesstrue \fi
}

\def\rref@linkpr#1{%
  \begingroup
  \reg@checkanchorless
  \edef\reg@refurl{\rref@current@url}%
  \ifreg@anchorless
    % No anchor by that name in this output, so no link.
    \endgroup#1%
  \else
    \ifx\reg@refurl\@empty
      \endgroup\href{\rref@current@fullurl}{#1}%
    \else
      % Another document, not published as HTML, so no link.
      \endgroup#1%
    \fi
  \fi
}
%    \end{macrocode}
% \iffalse
%</html>
% \fi
%
% \subsubsection{A known limitation: references to an item of \texttt{paras}}
%
%
% References to an article work; references to a paragraph do not. A legal text cites
% both -- `het in eerste lid genoemde bedrag' is as common as `zie artikel 3' --
% so this is worth knowing about rather than discovering.
%
% The cause is a name clash, not a missing anchor. For a list item tex4ht emits
% its own generic anchor (\texttt{x1-31}), while hyperref's \cmd{\refstepcounter} has
% meanwhile recorded \verb|\@currentHref| as \texttt{parasi.1.1}. regulatory stores that
% recorded value in zref, and \cmd{\aref} links to it -- so the link points at a name
% nothing carries. Verified on a two-item probe:
%
% \begin{center}
%     \begin{tabular}{@{}ll@{}}
%         ids in the output & \texttt{x1-2r1}, \texttt{article.1}, \texttt{x1-31} \\
%         link target       & \texttt{parasi.1.1} \\
%     \end{tabular}
% \end{center}
%
% What was tried and does \emph{not} work: hooking \cmd{\refstepcounter} to emit an extra
% anchor with \verb|\@currentHref|, both directly and deferred to \cmd{\AtBeginDocument}.
% Neither errors and neither produces an anchor -- \cmd{\refstepcounter} runs inside
% the construction of the item's label, and output written there is discarded.
% A working fix has to hook the item itself, not the counter.
%
% Left unfixed deliberately rather than half-fixed: an approach that silently
% does nothing is worse in this file than an absence that is written down.
% How a package should attach a hyperref anchor to an enumitem item is a question
% for the tex4ht maintainer.
%
% \subsubsection{Attachments}
%
%
% \cmd{\attachdocument} and \cmd{\documentattachment} do two things: embed a file inside the
% PDF, and link to that embedded copy. Neither half survives this output format,
% and both fail as an error rather than as an absence.
%
% \cmd{\embedfile} counts what it has embedded in \verb|\g__pdfmanagement_EmbeddedFiles_int|,
% which exists only while pdfmanagement is active -- it is not, here, because
% there is no \cmd{\DocumentMetadata} in an HTML run. The link is built from
% \verb|\pdfannot_link_begin:nnw|, an l3pdfannot command with no meaning outside a PDF
% backend. A single footnote carrying an attachment therefore produced six
% `Undefined control sequence' errors and make4ht stopped.
%
% An HTML page has no attachments and cannot acquire any, so there is nothing to
% translate the construct into. Both halves become no-ops and the link text is
% typeset as text -- the same principle the cross-document references follow
% above: a link this format cannot honour must not look like a link.
%
% Where that text argument happens to be a public URL the reader still sees where
% to go. Do not turn it into an \cmd{\href} on the strength of such a case:
% \cmd{\documentattachment} passes a document name in the same position, and a blanket
% hyperlink would point that at nothing.
%
% The replacement is made at the two public commands, not at \cmd{\embedfile} and
% \verb|\regulatory@attachmentlink| underneath them. Neutralising those two looks
% tighter and does not hold: \cmd{\embedfile} is set up again by its own package at
% begin-document, so a \cmd{\let} placed in a hook here is overwritten and the errors
% come back -- measured, four of them left. Redefining the public commands
% removes the only paths that reach \cmd{\embedfile} at all (\verb|\regulatory@embedonce| is
% called from \cmd{\documentattachment} and nowhere else), so nothing downstream can
% put it back.
%
% Parameters are taken outside the hook and the hook only \cmd{\let}s. Writing
% \cmd{\renewcommand}...{\#2} inside a hook body does not survive the round of
% parameter reduction the token list passes through: \#2 arrives as a literal
% digit, so the footnote reads its own parameter number instead of the text.
% \iffalse
%<*html>
% \fi
%    \begin{macrocode}
\newcommand\reg@attachdocument@text[3][]{#3}
\newcommand\reg@documentattachment@text[2]{#2}
\AtBeginDocument{%
  \let\attachdocument\reg@attachdocument@text
  \let\documentattachment\reg@documentattachment@text
}
%    \end{macrocode}
% \iffalse
%</html>
% \fi
%
% \subsubsection{The definitions list}
%
%
% The article that defines a contract's terms came out as one flat paragraph:
% every term and its description separated by nothing but whitespace. For the
% one article whose job is to distinguish a term from its meaning, that is the
% worst place in the document to lose the distinction.
%
% It comes from KOMA-Script's \texttt{labeling}, which the legal documents use to list
% the definitions explicitly:
%
% \begin{quote}\ttfamily
% \textbackslash{}begin\{labeling\}\{Overeenkomst van opdracht\}\\
% \hspace*{2em}\textbackslash{}item[Algemene Voorwaarden] \textbackslash{}describe\{terms\}\\
% \hspace*{2em}\textbackslash{}item[Consument] \textbackslash{}describe\{consumer\}
% \end{quote}
%
% tex4ht configures \texttt{itemize}, \texttt{enumerate} and \texttt{description}, but nothing ships
% a configuration for \texttt{labeling} -- so it falls through to the generic handling
% and the item structure is lost.
%
% Worth recording because it is where the search naturally goes first: the
% glossary style is not involved. \package{regulatory} prints with \texttt{style=alttree}, and
% changing that has no effect at all here -- verified by redefining the style
% and by forcing \texttt{style=altlist} onto the print command. Neither changed
% anything, because these documents do not print a glossary at this point; they
% enumerate the terms by hand.
%
% A list of terms with their meanings is a description list, so \texttt{labeling} is
% configured as one. Modelled on tex4ht's own \texttt{description} configuration in
% docbook.4ht: five arguments -- open list, close list, start item, after label
% -- with \cmd{\end:itm} carrying the close tag of the previous item, because a \texttt{<dd>}
% can only be closed once the next \texttt{<dt>} arrives or the list ends.
% \iffalse
%<*html>
% \fi
%    \begin{macrocode}
\ConfigureList{labeling}%
   {\EndP\HCode{<dl class="regulatory-definitions">}%
      \PushMacro\end:itm
      \global\let\end:itm=\empty}
   {\PopMacro\end:itm \global\let\end:itm \end:itm
    \EndP\HCode{</dd></dl>}\ShowPar}
   {\end:itm \global\def\end:itm{\EndP\HCode{</dd>}}%
    \HCode{<dt class="regulatory-term">}}
   {\EndP\HCode{</dt>}\HCode{<dd class="regulatory-definition">}\par\ShowPar}

\AtEndDocument{\reg@closearticle}
%    \end{macrocode}
% \iffalse
%</html>
% \fi
%
% \subsubsection{Default markup}
%
%
% \texttt{<section>} plus a real heading, not a \texttt{<div>} with a styled first line. The
% number sits in its own \texttt{<span>} so a stylesheet can set it apart, while staying
% inside the \texttt{<h2>} -- a screen reader should announce "Artikel 1. Wettelijke
% kaders" as one heading, because that is what it is.
%
% h2 rather than h1: the document title is the h1 and an article sits under it,
% which keeps the heading outline valid. Emitting headings at all is the whole
% point; emitting them at the wrong level would trade one accessibility problem
% for another.
% \iffalse
%<*html>
% \fi
%    \begin{macrocode}

\Configure{regarticle}
  {\HCode{<section class="regulatory-article"}}
  {\HCode{<h2 class="regulatory-article-heading"><span class="regulatory-article-number">}}
  {\HCode{</span> }}
  {\HCode{</h2>}}

\Configure{regpara}
  {\HCode{<section class="regulatory-para"}}
  {\HCode{<h3 class="regulatory-para-heading"><span class="regulatory-para-number">}}
  {\HCode{</span> }}
  {\HCode{</h3>}}
%    \end{macrocode}
% \iffalse
%</html>
% \fi
%
% \subsubsection{Translations}
%
%
% Deferred for the same reason as the redefinitions: regulatory.sty requires the
% translations package around line 99, well after this file is read. Calling
% \cmd{\providetranslation} here makes it an undefined control sequence, its braced
% arguments are typeset as text, and that ends the preamble.
% \iffalse
%<*html>
% \fi
%    \begin{macrocode}

\AtBeginDocument{%
  \providetranslation{Article}{Article}%
  \providetranslation[to=dutch]{Article}{Artikel}%
}

\endinput
%    \end{macrocode}
% \iffalse
%</html>
% \fi
%
% \Finale
%
